% !PRINTED 111
\weberpage{111}
%<<CONT
and further the sum

\weq{4}{x_0 + x_1 + x_2 + x_3 = 4A}

are rational numbers.

If we extract the square roots from (2) and take into account that, by (3),
$\sqrt{a}$ differs from $\sqrt{bc}$ only by a rational factor, then by
addition it follows that:

\weq{5}{x_0 = A + B\sqrt{b} + C\sqrt{c} + D\sqrt{bc},}

in which $A, B, C, D$ are rational numbers; and from this one obtains $x_1,
x_2, x_3$ by changing the signs of $\sqrt{b}, \sqrt{c}$.

But by Band I, \S.\,171, all square roots can, if necessary with the aid of
$i$, which is itself a cyclotomic number, be expressed rationally through
the cyclotomic periods (the \gesperrt{Gauss} sums), so that the proof of
our theorem is already contained in (5).

\noindent\hrulefill

If the group of the biquadratic equation is cyclic, we can arrange the
roots so that the group consists of the permutations

\weq{6}{1, \ \ (0,1,2,3), \ \ (0,2)(1,3), \ \ (0,3,2,1)}

and then we are to set:

\weq{7}{\begin{aligned}
4A &= x_0+x_1+x_2+x_3 \\
(i,x_0) &= x_0+ix_1-x_2-ix_3 \\
(-1,x_0) &= x_0-x_1+x_2-x_3 \\
(-i,x_0) &= x_0-ix_1-x_2+ix_3;
\end{aligned}}

in these, $A$ and $(-1,x_0)^2=m$ are rational, and hence $(-1,x_0)$ can be
expressed through cyclotomic numbers. The fourth powers

\weq{8}{(i,x_0)^4 = a+bi, \quad (-i,x_0)^4 = a-bi}

are numbers of the field $R(i)$, and in order to be able to draw the same
conclusions as with the cubic equations, it thus only remains to represent
$a+bi$ and $a-bi$ as fourth powers of cyclotomic numbers. For this there
serves further the theorem that

\weq{9}{(i,x_0)(-i,x_0) = c}

is a rational number.
